\newpage
\addcontentsline{toc}{section}{参考文献}
\begin{thebibliography}{99}
\bibitem{ref1} Mitra T, Gilbert E. The language that gets people to give: Phrases that predict success on Kickstarter[C]//Proceedings of the 17th ACM Conference on Computer Supported Cooperative Work \& Social Computing. 2014: 49-61.
\bibitem{ref2} Yuan H, Lau R Y K, Xu W. The determinants of crowdfunding success: A semantic text analytics approach[J]. Decision Support Systems, 2016, 91: 67-76.
\bibitem{ref3} Popescul D, Radu L D, P\u{a}v\u{a}loaia V D, et al. Psychological determinants of investor motivation in social media-based crowdfunding projects: A systematic review[J]. Frontiers in Psychology, 2020, 11: 588121.
\bibitem{ref4} Mollick E. The dynamics of crowdfunding: An exploratory study[J]. Journal of Business Venturing, 2014, 29(1): 1-16.
\bibitem{ref5} Pekar V, Candi M, Beltagui A, et al. Explainable text-based features in predictive models of crowdfunding campaigns[J]. Annals of Operations Research, 2025, 354(1): 367-397.
\bibitem{ref6} Ardakani S P, Hu J, Zhang J, et al. Identifying crowdfunding storytellers who deliver successful projects: A machine learning approach[J]. The Journal of Supercomputing, 2025, 81(1): 263.
\bibitem{ref7} Zhang X, Lyu H, Luo J. What contributes to a crowdfunding campaign's success? Evidence and analyses from GoFundMe data[J]. Journal of Social Computing, 2021, 2(2): 183-192.
\bibitem{ref8} Zhang X, Zheng X, Luo J. The influence of stereotypes and visual features on donation intentions in online medical crowdfunding campaigns: A comparison of survey and large language model-based methods[J]. npj Artificial Intelligence, 2025, 1(1): 16.
\bibitem{ref9} 张卫国, 黄思颖, 王超. 奖励众筹融资绩效动态预测研究——来自“众筹网”数据的实证[J]. 中国管理科学, 2023, 31(5): 71-83.
\bibitem{ref10} 张凯, 赵洪科, 刘淇, 等. 基于动静态表征的众筹协同预测方法[J]. 软件学报, 2020, 31(4): 967-980.
\bibitem{ref11} Greenberg M D, Pardo B, Hariharan K, et al. Crowdfunding support tools: Predicting success \& failure[C]//CHI '13 Extended Abstracts on Human Factors in Computing Systems. 2013: 1815-1820.
\bibitem{ref12} Li Y, Rakesh V, Reddy C K. Project success prediction in crowdfunding environments[C]//Proceedings of the Ninth ACM International Conference on Web Search and Data Mining. 2016: 247-256.
\bibitem{ref13} Kaminski J C, Hopp C. Predicting outcomes in crowdfunding campaigns with textual, visual, and linguistic signals[J]. Small Business Economics, 2020, 55(3): 627-649.
\bibitem{ref14} Cheng C, Tan F, Hou X, et al. Success prediction on crowdfunding with multimodal deep learning[C]//Proceedings of the 28th International Joint Conference on Artificial Intelligence. 2019: 2158-2164.
\bibitem{ref15} Tang Z, Yang Y, Li W, et al. Deep cross-attention network for crowdfunding success prediction[J]. IEEE Transactions on Multimedia, 2022, 25: 1306-1319.
\bibitem{ref16} Cai Z, Ding H, Xu M, et al. Multimodal dynamic graph convolutional network for crowdfunding success prediction[J]. Applied Soft Computing, 2024, 154: 111313.
\bibitem{ref17} Li J, Xu X, Li Y, et al. Commonality augmented disentanglement for multimodal crowdfunding success prediction[C]//Proceedings of the 2025 IEEE International Conference on Acoustics, Speech and Signal Processing. 2025: 1-5.
\bibitem{ref18} 黄健青, 黄晓凤, 殷国鹏. 众筹项目融资成功的影响因素及预测模型研究[J]. 中国软科学, 2017(7): 91-100.
\bibitem{ref19} 李小倩, 赵文红, 吕斯尧, 等. 产品价值线索对众筹融资绩效的影响——创业者线索的调节作用[J]. 研究与发展管理, 2024, 36(2): 125-138.
\bibitem{ref20} Kao T W, Hsiao S H, Su H C, et al. Deriving execution effectiveness of crowdfunding projects from the fundraiser network[J]. Journal of Management Information Systems, 2022, 39(1): 276-301.
\bibitem{ref21} 李清香, 王念新, 吕爽, 等. 发起人与出资者的在线交互对众筹项目成功的影响[J]. 管理工程学报, 2020, 34(1): 118-126.
\bibitem{ref22} 王念新, 吕爽, 周园, 等. 连续发起人的经验对众筹成功的影响: 经验相关性的调节效应分析[J]. 管理工程学报, 2020, 34(4): 89-100.
\bibitem{ref23} Carradini S, Fleischmann C. The effects of multimodal elements on success in Kickstarter crowdfunding campaigns[J]. Journal of Business and Technical Communication, 2023, 37(1): 1-27.
\bibitem{ref24} Raab M, Schlauderer S, Overhage S, et al. More than a feeling: Investigating the contagious effect of facial emotional expressions on investment decisions in reward-based crowdfunding[J]. Decision Support Systems, 2020, 135: 113326.
\bibitem{ref25} Costello F J, Lee K C. Exploring investors' expectancies and its impact on project funding success likelihood in crowdfunding by using text analytics and Bayesian networks[J]. Decision Support Systems, 2022, 154: 113695.
\bibitem{ref26} 刘畅, 陈守明. 文本信息含量与公益众筹绩效研究——基于有限注意的视角[J]. 投资研究, 2022, 41(6): 96-113.
\bibitem{ref27} Chen S, Wang H, Fang Y, et al. Informational and emotional appeals of cover image in crowdfunding platforms and the moderating role of campaign outputs[J]. Decision Support Systems, 2023, 171: 113975.
\bibitem{ref28} Song C, Luo J, H\"oltt\"a-Otto K, et al. Crowdfunding for design innovation: Prediction model with critical factors[J]. IEEE Transactions on Engineering Management, 2020, 69(4): 1565-1576.
\bibitem{ref29} Bao L, Wang Z, Zhao H. Who said what: Mining semantic features for success prediction in reward-based crowdfunding[J]. Electronic Commerce Research and Applications, 2022, 53: 101156.
\bibitem{ref30} Blanchard S J, Noseworthy T J, Pancer E, et al. Extraction of visual information to predict crowdfunding success[J]. Production and Operations Management, 2023, 32(12): 4172-4189.
\bibitem{ref31} 徐琳. 多模态数据驱动的众筹项目成功率预测研究[D]. 合肥工业大学, 2022.
\bibitem{ref32} Al-Qershi O M, Kwon J, Zhao S, et al. Predicting crowdfunding success with visuals and speech in video ads and text ads[J]. European Journal of Marketing, 2022, 56(6): 1610-1649.
\bibitem{ref33} Zhang Z, Lau R Y K. Exploiting multimodal features and deep learning for predicting crowdfunding successes[C]//Proceedings of the 2024 IEEE International Conference on Omni-layer Intelligent Systems. 2024: 1-6.
\bibitem{ref34} Sweller J. Cognitive load theory[M]//Psychology of Learning and Motivation. Academic Press, 2011, 55: 37-76.
\bibitem{ref35} 何旭, 郭春彦. 视觉工作记忆的容量与资源分配[J]. 心理科学进展, 2013, 21(10): 1741-1748.
\bibitem{ref36} 车敬上, 孙海龙, 肖晨洁, 等. 为什么信息超载损害决策? 基于有限认知资源的解释[J]. 心理科学进展, 2019, 27(10): 1758-1768.
\bibitem{ref37} Djamasbi S, Siegel M, Tullis T. Visual hierarchy and viewing behavior: An eye tracking study[C]//Human-Computer Interaction: Design and Development Approaches. 2011: 331-340.
\bibitem{ref38} Liang X, Hu X, Jiang J. Research on the effects of information description on crowdfunding success within a sustainable economy—the perspective of information communication[J]. Sustainability, 2020, 12(2): 650.
\bibitem{ref39} Alter A L, Oppenheimer D M. Uniting the tribes of fluency to form a metacognitive nation[J]. Personality and Social Psychology Review, 2009, 13(3): 219-235.
\bibitem{ref40} Hogarth R M, Einhorn H J. Order effects in belief updating: The belief-adjustment model[J]. Cognitive Psychology, 1992, 24(1): 1-55.
\bibitem{ref41} Joshi A, Mathur G. The inverted pyramid approach in user interface design for interactive information retrieval[C]//Proceedings of the EASY3 (CHI South India) Annual Conference. 2004.
\bibitem{ref42} Tjärnhage A, Söderström U, Norberg O, et al. The impact of scrollytelling on the reading experience of long-form journalism[C]//Proceedings of the European Conference on Cognitive Ergonomics 2023. 2023: 1-9.
\bibitem{ref43} Lee J E, Shin E, Kincade D H. The effects of presentation order of apparel product images on consumers' information processing style and purchase intentions[C]//Proceedings of the International Textile and Apparel Association Annual Conference. 2019, 76(1).
\bibitem{ref44} Lee J E, Shin E, Kincade D H. Presentation-order effect of product images on consumers' evaluations in online shopping[C]//Proceedings of the International Textile and Apparel Association Annual Conference. 2020, 77(1).
\bibitem{ref45} Razak N A, Asma'Amran S N. Effective language use in advertising in the most visible online stores: 2 Case Studies[J]. International Journal of Education and Learning Systems, 2017, 2: 30-38.
\bibitem{ref46} Green M C, Brock T C. The role of transportation in the persuasiveness of public narratives[J]. Journal of Personality and Social Psychology, 2000, 79(5): 701-721.
\bibitem{ref47} Xu Y, Li M, Cui L, et al. LayoutLM: Pre-training of text and layout for document image understanding[C]//Proceedings of the 26th ACM SIGKDD International Conference on Knowledge Discovery \& Data Mining. 2020: 1192-1200.
\bibitem{ref48} Xu Y, Xu Y, Lv T, et al. LayoutLMv2: Multi-modal pre-training for visually-rich document understanding[C]//Proceedings of the 59th Annual Meeting of the Association for Computational Linguistics and the 11th International Joint Conference on Natural Language Processing. 2021: 2579-2591.
\bibitem{ref49} Cui X, Lu W, Tong Y, et al. Diffusion-based multi-modal synergy interest network for click-through rate prediction[C]//Proceedings of the 48th International ACM SIGIR Conference on Research and Development in Information Retrieval. 2025: 581-591.
\bibitem{ref50} Chai S, Chang Y, Li Y. Image-left, text-right: The horizontal sequence effect of images and text on consumer responses to in-feed news[J]. Psychology \& Marketing, 2025, 42(8): 2121-2135.
\bibitem{ref51} Paivio A. Mental representations: A dual coding approach[M]. Oxford University Press, 1990.
\bibitem{ref52} Radford A, Kim J W, Hallacy C, et al. Learning transferable visual models from natural language supervision[C]//Proceedings of the 38th International Conference on Machine Learning. 2021: 8748-8763.
\end{thebibliography}
